% \documentclass{article}
% \usepackage{doc,url,verbatim,fancyvrb}
% \usepackage{pifont}
% \usepackage{pkgheads}
% \usepackage[letterpaper,body={6.3in,9.15in},top=.8in,left=1.1in]{geometry}
% \usepackage[hyperfootnotes=false]{hyperref}

\documentclass{article}
\usepackage[letterpaper,body={6.1in,9in},top=.8in,left=1.2in]{geometry}
\usepackage[no-math]{fontspec}
\usepackage{verbatim,fancyvrb}
\usepackage{gretl}
%\usepackage{amsmath}
%\usepackage{graphicx}
\usepackage[authoryear]{natbib}
\usepackage{hyperref}

\hypersetup{pdftitle={ordered\_predict 2.0},
            pdfauthor={Allin Cottrell},
            colorlinks=true,
            linkcolor=blue,
            urlcolor=red,
            citecolor=steel,
            bookmarksnumbered=true,
            plainpages=false
}

% \usepackage[utf8]{inputenc}
% \DeclareUnicodeCharacter{2212}{-}

\begin{document}

\setlength{\parindent}{0pt}
\setlength{\parskip}{1ex}
\setcounter{secnumdepth}{2}

\newcommand{\argname}[1]{\textsl{#1}}
  
\title{The \texttt{ordered\_predict} package, version 1.0}
\author{Allin Cottrell}
\maketitle

\section{Introduction}
\label{sec:intro}

This package computes outcome probabilities and predictions based on
results from an ordered response model (probit or logit). To set the
scene we begin with a brief review of such models. (See section 35.2
of the \textit{Gretl User's Guide} for more details.)

Let $y$ be an ordered response variable that can take on any of the
$J+1$ values ${0,1,2,\dots,J}$.  We suppose that underlying the
observed response is a latent variable,
\begin{equation}
  \label{eq:latent}
  y^* = X\beta + \varepsilon = z + \varepsilon
\end{equation}
Now define ``cut points'', $\alpha_1 < \alpha_2 < \cdots < \alpha_J$,
such that
%
\begin{equation}
  \label{eq:cuts}
  \begin{array}{ll}
    y = 0 & \textrm{if } \, y^* \leq \alpha_1 \\
    y = 1 & \textrm{if } \, \alpha_1 < y^* \leq \alpha_2 \\
    \vdots \\
    y = J & \textrm{if } \, y^* > \alpha_J \\
  \end{array}
\end{equation}
For example, if the response takes on three values there will be two
such cut points, $\alpha_1$ and $\alpha_2$. 

The probability that individual $i$ exhibits response $j$, conditional
on the characteristics $x_i$, is then given by
%
\begin{equation}
  \label{eq:ordered-probs}
  P(y_i = j \,|\, x_i) = \left\{
    \begin{array}{ll}
      P(y^* \leq \alpha_1 \,|\, x_i) = F(\alpha_1 - z_i) & \textrm{for }\, j = 0 \\
      P(\alpha_j < y^* \leq \alpha_{j+1} \,|\, x_i) = 
      F(\alpha_{j+1} - z_i) -  F(\alpha_j - z_i) & \textrm{for }\, 0 < j < J \\
      P(y^* > \alpha_J \,|\, x_i) = 1 -  F(\alpha_J - z_i) & \textrm{for }\, j = J 
    \end{array}
    \right.
\end{equation}
%
where $F(\cdot)$ denotes a cumulative distribution function: either
the normal CDF (in the case of probit) or the logistic CDF
(logit). The unknown parameters $\alpha_j$ are estimated jointly with
the $\beta$ terms via maximum likelihood.

In gretl the \texttt{probit} or \texttt{logit} command can be used to
estimate such a model. After estimation the accessor \dollar{yhat} can
be used to retrieve the estimates $X\hat{\beta} = \hat{z}$ (see
equation~\ref{eq:latent}). This package can then be used to compute
the per-outcome probabilities as shown in (\ref{eq:ordered-probs}),
with inputs $\hat{z}$ (labeled \texttt{zhat} below) and the vector of
cut-point estimates (labeled \texttt{cut} below). The latter can be
obtained by selecting the appropriate trailing rows of the
\dollar{coeff} vector: if the specification includes $k$ independent
variables then one wants
\begin{code}
matrix cut = $coeff[k+1:]
\end{code}
since \dollar{coeff} holds $\hat\beta$ followed by $\hat\alpha$.

\section{Usage via script}
\label{sec:script-usage}

Two main public functions are offered for scripting purposes; they
each take two required arguments, namely a cut-point vector and a
$\hat{z}$ series (plus optional trailing arguments as described
below).

Note that for out-of-sample prediction one can supply a $\hat z$ other
than that obtained via the \dollar{yhat} accessor. For example,
suppose one has 500 observations and one wishes to estimate a model on
the first 400 and obtain predictions for the remaining 100. Then
after ordered estimation one can do
\begin{code}
smpl 401 500
series zhat = lincomb(Xlist, $coeff[1:nelem(Xlist)])
\end{code}
% $
where \texttt{Xlist} is the list of regressors in the specification.

\subsection{Obtaining the matrix of probabilities}

To obtain an outcome-probability matrix use the following function:
\begin{code}
matrix ordered_Pmat (const matrix cut,
                     series zhat,
                     bool lgt[0])
\end{code}
It returns an $n \times m$ matrix, where $n$ is the number of
observations and $m$ is the number of distinct response values (that
is, the number of cut points plus one). Row $i$ of this matrix holds
the estimated probability of each of the $m$ responses at observation
$i$. Give a non-zero value for the third, optional argument to
select the logit variant, the default being probit.

\subsection{Obtaining point predictions}
\label{sec:point-pred}

In some cases one may wish to ask: which particular outcome should be
predicted for case $i$ (integer-valued $\hat{y}_i$) conditional on the
characteristics $x_i$ along with $\hat\beta$ and $\hat\alpha$? The
package provides means of answering that question, but it is not as
straightforward as it may seem.

The most ``natural'' way of arriving at a point prediction is,
perhaps, to choose the outcome with the greatest estimated
probability. We'll call this the \textbf{pmax} method; it is easily
implemented given the outcome-probability matrix, as in
\begin{code}
matrix Pij = ordered_Pmat(cut, zhat)
matrix pmax = imaxr(Pij) - 1
\end{code}
The built-in \texttt{imaxr} function returns a column vector holding
the row indices of the maximum values on each row of the input matrix,
so in this case an $n$-vector with elements in the range 1 to $J+1$
(using the notation of section~\ref{sec:intro}, and where $n$ is the
length of \texttt{zhat}). We subtract 1 to get the zero-based value of
the maximum-probability outcome---again, as per
section~\ref{sec:intro}, but see also section~\ref{sec:ymin} below.

However, there are (at least) two other ways of arriving at a point
prediction.
\begin{itemize}
\item Under the assumption that $E(\varepsilon|X) = 0$,
  $E(\hat{z})$ should equal $z$ and one could assign a prediction
  based on where $\hat{z}$ falls in relation to the estimated cut
  points. That is, predict $y_i = 0$ when $\hat{z}_i \leq \alpha_1$,
  $y_i = 1$ when $\alpha_1 < \hat{z}_i \leq \alpha_2$, and so on---see
  (\ref{eq:cuts}). Call this the \textbf{cuts} method.
\item Use all the probability information for case $i$ by forming an
  expected value, as in
  \[
    E(y_i|x_i) = \sum_{j=0}^J \hat{P}(y_i = j \,|\, x_i)\cdot j
  \]
  which one could then round to the nearest integer. Call this the
  \textbf{expectation} method.
\end{itemize}

These three methods may well yield the same specific prediction in
many cases but they don't have to; they are not logically equivalent
and in fact can diverge quite substantially, with the extent of
the divergence depending on, among other things, the variance of the
error term, $\varepsilon$. So is one of the methods clearly better
than the others? It depends on the criterion for success in prediction
(or its flipside, the loss function). Here are our findings in brief
(see the Appendix for more details):

\begin{itemize}
\item To maximize the number of correct predictions, use the
  \textbf{pmax} method.
\item To minimize the sum of absolute prediction errors, use the
  \textbf{cuts} method.  
\item To minimize the number of cases where the prediction differs
  from the observed value by more than one step, use the
  \textbf{expectation} method.
\end{itemize}

All that said, here is the signature of the function for obtaining
point predictions.

\begin{code}
series ordered_prediction (const matrix cut,
                           series zhat,
                           int ymin[0],
                           bool lgt[0],
                           int ptype[1:3:1])
\end{code}
The return value is a series containing a predicted response at each
observation based on \texttt{cut} and \texttt{zhat}. The optional
arguments have the following effects:
\begin{itemize}
\item \texttt{ymin} can be used to adjust the minimum value of the
  response variable. By default this is taken to be zero, but if (for
  example) you give \texttt{ymin} = 3 the responses are labeled 3, 4,
  5,\dots\
\item \texttt{lgt} can be used to specify that the logit estimator
  should be assumed (the default being probit).
\item \texttt{ptype} can be used to specify a prediction method (1 =
  \textbf{pmax}, the default; 2 = \textbf{cuts}; 3 =
  \textbf{expectation}).
\end{itemize}

In addition, for purposes of comparison, the package offers a variant
on \texttt{ordered\_prediction()} which returns a matrix containing
predictions from all three methods, in the order indicated above. The
signature of this function is as follows,
\begin{code}
matrix all_ordered_predictions (const matrix cut,
                                series zhat,
                                int ymin[0],
                                bool lgt[0])
\end{code}
and the returned matrix is $n \times 3$, where $n$ is the number of
values of \texttt{zhat} in the current sample range.

\subsection{More on \texttt{ymin}}
\label{sec:ymin}

It may be helpful to further explain the \texttt{ymin} argument to
\texttt{ordered\_prediction}. In the theoretical presentation of the
ordered-outcome model in section~\ref{sec:intro} we took the minimum
value of $y$ to be zero, by construction, and this is also what gretl
does internally. In some applications, however, the actual $y$ data
may have a positive minimum. For example, suppose our dependent
variable is a rating on a scale of 0 to 10 (in principle) but in our
sample there are no values below 3 and none higher than 8. In effect,
therefore, we have six levels of $y$ (the integers 3 to 8) which will
be recoded internally as 0 to 5.

If we were to generate predictions of $y$ on the 0 to 5 scale and
check them for equality with observed $y$ to find the number of
correct predictions, we'd likely get a spurious answer of none
correct. For comparability we either have to subtract 3 from observed
$y$ or add 3 to the predictions. The \texttt{ymin} argument lets you
do the latter.

\section{GUI access}

On installing the \texttt{ordered\_predict} package you should get the
choice of letting it attach to the \textsf{Analysis} menu in gretl
model windows under the item \textsf{Outcome probabilities}. In that
case when you estimate an ordered model via the graphical interface
you can view the matrix of per-outcome probabilities by clicking on
this item. The output also identifies the outcome with the greatest
individual probability.

\clearpage

\section*{Appendix}

We present here some background for the statements made above
concerning the relative merits of three prediction methods for ordered
models (see section~\ref{sec:point-pred}).

We ran several Monte Carlo exercises (differing with respect to the
variance of the error term and the relative frequency of occurrence of
the discrete $y$ values) and obtained similar results on all of them,
at least with regard to the ordering of the three prediction methods
on each of the four criteria considered, which were:
\begin{enumerate}
\item Proportion of cases correctly predicted.
\item Proportion of cases where the prediction misses actual by more
  than one step.
\item Mean Absolute Error (MAE) of the prediction.
\item Mean Squared Error (MSE) of the prediction.
\end{enumerate}

With a dependent variable that is not on an interval scale one might
have doubts about the applicability of MAE and MSE (although it's
presumably ``worse'' even on an ordinal scale for the prediction to
miss the actual value by more steps); that's why we included the
relatively robust criterion 2.

The table below shows typical results for a simulation with a single
normally distributed independent variable and a normally distributed
error, where $y$ is discretized into 4 levels, with a sample size of
200 at each iteration and 10,000 iterations. The error variance was
calibrated to give disagreement between the methods in about 5 percent
of cases. The figures of merit in the table are based solely on the
(102,559) cases where the three methods were not unanimous, and they
are normalized to a value of 1.0 for best performance on each row
(that is, the greatest value for ``Correct predictions'' and the least
value for the other criteria).

\begin{center}
  \begin{tabular}{lrrr}
    & \multicolumn{3}{c}{Method} \\
    & \multicolumn{1}{c}{1} & \multicolumn{1}{c}{2} &
      \multicolumn{1}{c}{3} \\
    Correct predictions & 1.000 & 0.984 & 0.909 \\
    Cases off by more than one & 6.802 & 3.995 & 1.000 \\
    Mean Absolute Error & 1.024 & 1.000 & 1.027 \\
    Mean Squared Error & 1.147 & 1.052 & 1.000
  \end{tabular}

  {\small Methods: 1 = \textbf{pmax}, 2 = \textbf{cuts}, 3 =
  \textbf{expectation}}
\end{center}

Not surprisingly, when the simulation is calibrated to produce fewer
cases of disagreement the results tend to become more uniform (e.g.\
the counts of correct predictions for methods 2 and 3 get closer to
that of method 1). But in all our simulations the rank-order of the
methods on each criterion was invariant.

So it's noteworthy that while the \textbf{pmax} method tends to give
the greatest number of correct predictions, it also tends to give the
greatest number of substantially wrong predictions (off by more than
one). Seriously wrong predictions are minimized by the
\textbf{expectation} method, while the \textbf{cuts} method seems to
occupy the middle ground, with the virtue (if it be such) of tending
to minimize the MAE.

The hansl script which generated the above results is shown on the
following page. Note that with a large number of Monte Carlo
iterations some defensive programming is required: we need to handle
cases where the estimator fails to converge, and we also exclude cases
datasets where the dependent variable takes on fewer than its maximum
of 4 levels.

\clearpage

\begin{script}[htbp]
  \caption{Monte Carlo script for comparing prediction methods}
  \label{list:mc}
\begin{code}
set verbose off
include ordered_predict.gfn

nulldata 200
set seed 1544461131 # or change this

# artificial independent (latent) variable
series x = normal()

iters = 10000
K = iters * $nobs
scalar N = 0
scalar m = 0

# allow sloppier than usual BFGS convergence criterion
set bfgs_maxgrad 8.0

# 0.2 * K should be more than enough rows
matrix R = zeros(0.2*K, 3)

loop while N < K --quiet
    series z = x + normal()/1.6
    series y = z < -1.1 ? 0 : z < 0 ? 1 : z < 1.1 ? 2 : 3
    catch probit y x --quiet
    if $error
        nc = 0
    else
        c = $coeff[2:]
        nc = nelem(c)
    endif
    if nc == 3
        # successful estimation with 3 cuts
        N += $T
        matrix P = all_ordered_predictions(c, $yhat, 0, 0)
        matrix sel = !(P[,1] .= P[,2] && P[,2] .= P[,3])
        matrix d = selifr(abs(y .- P), sel)
        nd = rows(d)
        if nd > 0
            R[m+1:m+nd,] = d
            m += nd
        endif
    endif
endloop

R = R[1:m,] # restrict R to the m non-unanimous cases
printf "\nDiffering predictions: %d out of %d (%.1f%%)\n", m, N, 100*m/N
printf "\nIn cases of differing predictions --\n"
r = sumc(R .= 0)
printf "count of correct predictions:\n%8.3f", r * 1/r[imaxr(r)]
r = sumc(R .> 1)
printf "count of predictions off by more than one:\n%8.3f", r * 1/r[iminr(r)]
r = sumc(R)
printf "sum of absolute errors:\n%8.3f", r * 1/r[iminr(r)]
r = sumc(R.^2)
printf "sum of squared errors:\n%8.3f", r * 1/r[iminr(r)]
\end{code}
\end{script}


\end{document}
